\hypertarget{05-engine-i-ch024}{}
\chapter{Profiling and language/platform ownership}\label{stagechapter:05-engine-i-ch024}\label{05-engine-i-ch024-profiling-and-languageplatform-ownership}

\hyperref[volume-reading-map]{Reading map} · \hyperref[05-engine-i-ch023-execute-and-measure-a-metal-frame]{Previous chapter} · \hyperref[06-engine-ii-ch025-renderer-evolution-through-resource-ownership]{Next chapter}

A renderer spans profiling tools, Objective-C++, Swift, and platform state, each with different ownership rules. This chapter keeps measurements scoped to the right build and makes language and process-wide state restoration explicit.

\textbf{By the end:} Interpret profiling output with build and workload context, and cross language/platform boundaries without leaking ownership or global state.

\textbf{Recall question:} What does this profile actually measure, and what state must be restored when the platform-specific operation ends?

\subsection{PROFILING, OBJECTIVE-C++ OWNERSHIP, SWIFTUI \& THREAD BOUNDARIES}\label{05-engine-i-ch024-profiling-objective-c-ownership-swiftui-thread-boundaries}

\textbf{GPU counters and Metal profiling}

You should distinguish at least three kinds of performance measurements.

CPU encoding time:

\[
\text{how long CPU spent building commands}.
\]

Command-buffer GPU duration:

\[
\text{how long GPU command buffer occupied GPU timeline}.
\]

Detailed GPU counters:

\[
\text{what hardware was doing inside that duration}.
\]

Those answer different questions.

\begin{center}\rule{0.5\linewidth}{0.5pt}\end{center}

\textbf{Maveb's current coarse GPU timing}

Completion handler reads:

\begin{verbatim}
completed->GPUEndTime()
-
completed->GPUStartTime()
\end{verbatim}

then converts to milliseconds.

This gives whole-command-buffer GPU duration.

Useful for:

\[
\text{frame GPU ms}.
\]

But not sufficient to tell whether:

\begin{verbatim}
Gaussian sort
shadow
TAA
bloom
presentation
\end{verbatim}

is slow.

\begin{center}\rule{0.5\linewidth}{0.5pt}\end{center}

\textbf{Counter sample buffers}

Modern Metal profiling can use counter sampling mechanisms where supported.

Conceptually record counter samples at boundaries:

\begin{verbatim}
sample 0
Gaussian projection
sample 1
Radix
sample 2
Composite
sample 3
\end{verbatim}

Then resolve/read results after GPU completion.

Counters can expose architecture-dependent metrics such as timestamps and other supported statistics.

Capabilities must be queried.

\begin{center}\rule{0.5\linewidth}{0.5pt}\end{center}

\textbf{Profiling should not stall}

Bad profiler:

\begin{verbatim}
commit
waitUntilCompleted
read counters
\end{verbatim}

every frame.

That destroys CPU/GPU overlap.

Better:

\begin{verbatim}
frame N records counters
frame N+2 reads frame N when already completed
\end{verbatim}

This is the same latency-hiding principle as Vulkan query pools.

\begin{center}\rule{0.5\linewidth}{0.5pt}\end{center}

\textbf{Xcode GPU capture}

Xcode's Metal GPU capture can inspect:

\begin{itemize}
\tightlist
\item
  command buffers,
\item
  encoders,
\item
  resources,
\item
  pipeline state,
\item
  textures,
\item
  draw calls,
\item
  shader execution,
\item
  performance.
\end{itemize}

This is the Metal equivalent of the workflow you might use RenderDoc for Vulkan, though Apple's integration can expose Apple-specific GPU details.

\begin{center}\rule{0.5\linewidth}{0.5pt}\end{center}

\textbf{Capture tells you execution reality}

Your source might say:

\begin{verbatim}
Render scene
\end{verbatim}

but GPU capture may reveal:

\begin{itemize}
\tightlist
\item
  five unexpected render passes,
\item
  attachment stored and reloaded repeatedly,
\item
  gigantic intermediate texture,
\item
  redundant state binding,
\item
  unexpected shader variant.
\end{itemize}

Performance optimization should eventually inspect the executed GPU workload, not merely source code.

\begin{center}\rule{0.5\linewidth}{0.5pt}\end{center}

\textbf{Metal labels}

Metal resources and command objects expose labels.

Maveb labels extensively:

\begin{verbatim}
AETHER Primary Queue
AETHER Frame
Scene HDR Color
Scene Reverse-Z Depth
Gaussian HDR Color
Temporal Resolve Pass
Presentation Tone Map Pass
Mesh Upload
Presented Frame Capture
\end{verbatim}

This is excellent.

It means GPU capture shows semantic engine names rather than anonymous resources.

\begin{center}\rule{0.5\linewidth}{0.5pt}\end{center}

\textbf{Encoder labels}

For every meaningful pass:

\begin{verbatim}
encoder->
    setLabel(
        "Gaussian Radix Histogram");
\end{verbatim}

or equivalent.

You can also use debug groups where useful:

\begin{verbatim}
Frame
    Shadow
    Scene
    Gaussian
        Project
        Sort
        Composite
    TAA
    Bloom
    Present
\end{verbatim}

A profiler/capture should reflect your engine's conceptual hierarchy.

\begin{center}\rule{0.5\linewidth}{0.5pt}\end{center}

\textbf{Why labels are not optional fluff}

When debugging a 200-pass renderer, names save enormous time.

Compare:

\begin{verbatim}
Texture 0x173902
Encoder 0x77193
\end{verbatim}

with:

\begin{verbatim}
Scene Motion Vectors
Temporal Resolve Pass
\end{verbatim}

Labels are part of engine observability.

\begin{center}\rule{0.5\linewidth}{0.5pt}\end{center}

\textbf{Objective-C retain/release intuition under ARC}

Objective-C objects use reference counting.

Conceptually:

\[
retain
\Rightarrow
refcount+1,
\]

\[
release
\Rightarrow
refcount-1.
\]

When count reaches zero:

\[
object\ destroyed.
\]

ARC automatically inserts many retain/release operations based on ownership semantics.

\begin{center}\rule{0.5\linewidth}{0.5pt}\end{center}

\textbf{Strong ownership}

A strong reference keeps object alive.

Objective-C property:

\begin{verbatim}
@property(strong)
id<MTLDevice> device;
\end{verbatim}

means owner retains it.

When owner releases property, reference count can fall.

This is analogous conceptually to \texttt{shared\_ptr}, but implementations and semantics are not identical.

\begin{center}\rule{0.5\linewidth}{0.5pt}\end{center}

\textbf{Weak references}

Weak reference:

\[
\boxed{\text{does not keep object alive}}
\]

and is automatically zeroed when referenced Objective-C object dies.

Useful for delegate relationships to avoid retain cycles.

Again, this is Objective-C object lifetime, not GPU completion lifetime.

A Metal object can be strongly alive while GPU work is unfinished, and command-buffer/resource retention semantics matter too.

\begin{center}\rule{0.5\linewidth}{0.5pt}\end{center}

\textbf{Autorelease pools}

Some Objective-C APIs return autoreleased temporary objects.

A tight render loop without a local autorelease pool can accumulate temporary objects until outer pool drains.

Therefore render code frequently uses:

\begin{verbatim}
@autoreleasepool
{
    ...
}
\end{verbatim}

Kairo does this in Metal draw calls.

Maveb's C++ Metal renderer explicitly creates/releases an autorelease pool per frame.

That is correct.

\begin{center}\rule{0.5\linewidth}{0.5pt}\end{center}

\textbf{ARC does not manage C++ objects}

Suppose Objective-C++ object contains:

\begin{verbatim}
std::unique_ptr<
    aether::metal::Renderer>
    _renderer;
\end{verbatim}

ARC does not manage that pointer.

C++ RAII does.

So in Objective-C++ source you have \textbf{two ownership systems simultaneously}:

\[
\boxed{
\text{ARC for Objective-C objects}
}
\]

and:

\[
\boxed{
\text{C++ RAII for C++ objects}.
}
\]

You must understand which system owns each object.

\begin{center}\rule{0.5\linewidth}{0.5pt}\end{center}

\textbf{metal-cpp adds a third-looking syntax, but same ownership reality}

Maveb uses C++ wrappers around Objective-C Metal objects.

Its helper pattern distinguishes:

\[
adopt(x)
\]

from:

\[
retain(x).
\]

Conceptually:

\texttt{adopt()} means:

\begin{quote}
this newly created object already carries ownership; wrap it without adding another retain.
\end{quote}

\texttt{retain()} means:

\begin{quote}
borrowed object must be retained because wrapper will own it.
\end{quote}

Mix those incorrectly and you get:

\begin{itemize}
\tightlist
\item
  leaks,
\item
  double releases,
\item
  dangling pointers.
\end{itemize}

\begin{center}\rule{0.5\linewidth}{0.5pt}\end{center}

\textbf{Objective-C++ \texttt{.mm} translation units}

A \texttt{.mm} file can contain both:

\begin{verbatim}
id<MTLDevice>
\end{verbatim}

and:

\begin{verbatim}
std::unique_ptr<Renderer>
\end{verbatim}

and Objective-C message syntax:

\begin{verbatim}
[object method]
\end{verbatim}

plus ordinary modern C++.

This is the ideal bridge layer between:

\[
\text{C++ engine core}
\]

and:

\[
\text{AppKit/Metal/Objective-C APIs}.
\]

\begin{center}\rule{0.5\linewidth}{0.5pt}\end{center}

\textbf{Why not put Objective-C in all C++ files?}

It infects portability.

If core engine headers contain:

\begin{verbatim}
#import <Metal/Metal.h>
\end{verbatim}

or Objective-C object syntax, Linux/Windows compilers cannot parse them.

Better architecture:

\begin{verbatim}
portable C++ headers
        ↓
Objective-C++ implementation bridge (.mm)
        ↓
Metal / AppKit
\end{verbatim}

This keeps platform code localized.

\begin{center}\rule{0.5\linewidth}{0.5pt}\end{center}

\textbf{Kairo follows this idea}

Public module:

\begin{verbatim}
MetalRenderer.cppm
\end{verbatim}

talks through:

\begin{verbatim}
detail::MetalBackend
\end{verbatim}

while raw Objective-C Metal implementation lives in:

\begin{verbatim}
detail/MetalBackend.mm.
\end{verbatim}

That is a good boundary.

The public renderer-neutral host does not need Objective-C syntax leaking through every module.

\begin{center}\rule{0.5\linewidth}{0.5pt}\end{center}

\textbf{C++ ownership across Objective-C++ bridges}

AetherStudio's bridge contains:

\begin{verbatim}
std::unique_ptr<
    aether::metal::Renderer>
    _renderer;
\end{verbatim}

inside:

\begin{verbatim}
@implementation
AetherViewportDelegate
\end{verbatim}

This is a very clean boundary.

Objective-C object owns a C++ unique pointer.

When delegate dies, C++ renderer destructor executes.

Renderer then releases its own Metal C++ wrapper resources through its RAII mechanisms.

\begin{center}\rule{0.5\linewidth}{0.5pt}\end{center}

\textbf{Passing an Objective-C Metal device into C++}

Aether bridge:

\begin{verbatim}
id<MTLDevice> device
\end{verbatim}

is converted to metal-cpp representation with a non-owning bridge cast.

Then:

\begin{verbatim}
Renderer::create(
    reinterpret_cast<
        MTL::Device*>(
        (__bridge void*)device));
\end{verbatim}

The \texttt{\_\_bridge} cast does not transfer Objective-C ownership.

Therefore C++ renderer must establish its own ownership if it stores the device.

Maveb constructor does:

\begin{verbatim}
device_(retain(device))
\end{verbatim}

which is exactly what you want.

\begin{center}\rule{0.5\linewidth}{0.5pt}\end{center}

\textbf{Ownership equation}

At bridge:

\[
\boxed{
\text{borrow ObjC device}
}
\]

then inside persistent C++ renderer:

\[
\boxed{
\text{retain device for C++ lifetime}.
}
\]

If it merely stored borrowed pointer, Objective-C owner could disappear and leave dangling C++ Metal pointer.

The retain is the lifetime bridge.

\begin{center}\rule{0.5\linewidth}{0.5pt}\end{center}

\textbf{error bridging C++ → Objective-C/Swift}

Never casually allow C++ exceptions to cross into Swift/Objective-C API boundaries.

Different exception runtimes/ABI semantics make that unsafe architecture.

Instead transform:

\[
C++\ Result<T,Error>
\]

into:

\[
Objective\text{-}C\ return\ value
+
NSError.
\]

\begin{center}\rule{0.5\linewidth}{0.5pt}\end{center}

\textbf{Example bridge}

C++:

\begin{verbatim}
Result<void>
SaveThing(...);
\end{verbatim}

Objective-C++:

\begin{verbatim}
BOOL SaveThingBridge(
    NSURL* destination,
    NSError** error)
{
    auto result =
        SaveThing(...);

    if (!result)
    {
        if (error)
        {
            *error =
                [NSError
                    errorWithDomain:
                        @"com.engine"
                    code:
                        1
                    userInfo:
                        @{
                          NSLocalizedDescriptionKey:
                            ...
                        }];
        }

        return NO;
    }

    return YES;
}
\end{verbatim}

This is exactly the pattern Aether uses in \texttt{AetherWriteDiagnostics}.

\begin{center}\rule{0.5\linewidth}{0.5pt}\end{center}

\textbf{Current Aether error bridge has multiple levels}

\texttt{AetherWriteDiagnostics()} converts failures into:

\[
NSError**.
\]

Many viewport operations currently instead:

\begin{itemize}
\tightlist
\item
  return \texttt{BOOL},
\item
  return \texttt{nil},
\item
  return empty arrays,
\item
  log engine error.
\end{itemize}

That is acceptable for internal editor operations.

For APIs where caller must distinguish failure reasons, promote them to structured \texttt{NSError}/Swift \texttt{Error}.

\begin{center}\rule{0.5\linewidth}{0.5pt}\end{center}

\textbf{Preserve error semantics}

Do not reduce:

\begin{verbatim}
resource exhausted
invalid argument
file not found
corrupt data
Metal command failure
\end{verbatim}

into:

\begin{verbatim}
"something failed"
\end{verbatim}

A bridge should preserve:

\[
error\ code,
\]

\[
message,
\]

\[
context.
\]

Maveb's internal errors already carry:

\[
ErrorCode
+
description/context.
\]

That is useful information to carry upward.

\begin{center}\rule{0.5\linewidth}{0.5pt}\end{center}

\textbf{SwiftUI view versus rendering-core responsibility}

This separation is extremely important.

SwiftUI should own:

\[
\text{application/UI state}.
\]

Renderer should own:

\[
\text{GPU resources and frame execution}.
\]

Bridge translates between them.

Do not let SwiftUI understand:

\begin{verbatim}
MTLBuffer
MTLTexture
command encoders
Gaussian radix buffers
BVH nodes
\end{verbatim}

and do not make render core understand:

\begin{verbatim}
@Binding
@State
@AppStorage
NSViewRepresentable
\end{verbatim}

\begin{center}\rule{0.5\linewidth}{0.5pt}\end{center}

\textbf{Aether's architecture}

SwiftUI:

\begin{verbatim}
AetherViewport
\end{verbatim}

contains bindings:

\begin{verbatim}
selectedGaussianId
selectedMeshId
camera
materials
lights
playback
debug settings
\end{verbatim}

It creates:

\begin{verbatim}
AetherViewportView
\end{verbatim}

through \texttt{NSViewRepresentable}.

Objective-C++ view owns:

\begin{verbatim}
MTKView
AetherViewportDelegate
\end{verbatim}

Delegate owns:

\begin{verbatim}
std::unique_ptr<
    aether::metal::Renderer>
\end{verbatim}

This is a strong layered architecture.

\begin{center}\rule{0.5\linewidth}{0.5pt}\end{center}

\textbf{The responsibility chain}

\begin{verbatim}
SwiftUI
    declarative editor state

        ↓

NSViewRepresentable
    lifecycle/update synchronization

        ↓

AetherViewportView
    AppKit input/window/view integration

        ↓

AetherViewportDelegate
    Objective-C++ conversion boundary

        ↓

aether::metal::Renderer
    C++ rendering system

        ↓

Metal
\end{verbatim}

Every layer speaks its natural language.

That is what good platform integration looks like.

\begin{center}\rule{0.5\linewidth}{0.5pt}\end{center}

\textbf{Why SwiftUI should not issue GPU commands}

SwiftUI body/update execution is driven by declarative state changes.

It is not a stable real-time rendering loop.

GPU frame execution needs:

\begin{itemize}
\tightlist
\item
  consistent timing,
\item
  resource lifetimes,
\item
  command-buffer scheduling,
\item
  frames in flight.
\end{itemize}

\texttt{MTKView}/renderer handles that.

SwiftUI configures it.

\begin{center}\rule{0.5\linewidth}{0.5pt}\end{center}

\textbf{\texttt{NSViewRepresentable}}

Aether:

\begin{verbatim}
struct AetherViewport:
    NSViewRepresentable
\end{verbatim}

implements:

\begin{verbatim}
makeNSView(...)
\end{verbatim}

and:

\begin{verbatim}
updateNSView(...)
\end{verbatim}

\texttt{makeNSView} constructs native Metal-backed view once.

\texttt{updateNSView} propagates changes in Swift state into existing native renderer/view.

That is exactly how a long-lived renderer should integrate with SwiftUI.

\begin{center}\rule{0.5\linewidth}{0.5pt}\end{center}

\textbf{Coordinator prevents redundant work}

Aether's \texttt{Coordinator} caches previously applied:

\begin{verbatim}
scenePath
dynamicMeshPath
transform overrides
materials
lights
camera
playback
\end{verbatim}

Then \texttt{updateNSView} compares.

Only changed values are pushed into renderer.

This matters because SwiftUI may call update many times.

You do not want to reload a scene or re-upload every material because some unrelated UI state changed.

\begin{center}\rule{0.5\linewidth}{0.5pt}\end{center}

\textbf{main-thread UI constraints}

AppKit and SwiftUI state/view operations generally belong on the main thread.

Examples:

\begin{verbatim}
creating/changing NSView hierarchy
mutating SwiftUI-observed UI state
handling many AppKit events
\end{verbatim}

GPU completion handlers, worker threads and asset loaders may run elsewhere.

The rule becomes:

\[
\boxed{
\text{renderer thread may produce data}
\rightarrow
\text{main thread consumes UI-facing result}.
}
\]

\begin{center}\rule{0.5\linewidth}{0.5pt}\end{center}

\textbf{Why this matters with Metal completion handlers}

Completion handler might do:

\begin{verbatim}
completedFrameCapture_ =
    capture;
\end{verbatim}

under mutex.

Good.

It should not directly do:

\begin{verbatim}
selectedImage = capture
\end{verbatim}

or manipulate \texttt{NSView} from unknown callback thread.

Instead:

\begin{verbatim}
GPU callback
    ↓
thread-safe result queue/state
    ↓
DispatchQueue.main
    ↓
UI update
\end{verbatim}

\begin{center}\rule{0.5\linewidth}{0.5pt}\end{center}

\textbf{Aether Swift code already uses main dispatch for deferred state updates}

Examples in \texttt{updateNSView}:

\begin{verbatim}
DispatchQueue.main.async {
    selectedMeshTransform =
        snapshot
}
\end{verbatim}

and similar material/selection updates.

This avoids mutating SwiftUI binding synchronously in problematic update cycles and keeps UI mutation on main.

\begin{center}\rule{0.5\linewidth}{0.5pt}\end{center}

\textbf{cross-thread renderer callbacks}

Once renderer becomes asynchronous, events can originate from:

\begin{verbatim}
GPU completion
asset loader
simulation
network
background reconstruction
\end{verbatim}

while UI exists on main thread.

A callback boundary therefore needs an explicit threading contract.

For every callback answer:

\[
\boxed{
\text{What thread invokes it?}
}
\]

\[
\boxed{
\text{Who owns captured data?}
}
\]

\[
\boxed{
\text{May receiver outlive sender?}
}
\]

\[
\boxed{
\text{Must callback be marshalled to main thread?}
}
\]

\begin{center}\rule{0.5\linewidth}{0.5pt}\end{center}

\textbf{Bad cross-thread callback}

Imagine:

\begin{verbatim}
commandBuffer->
    addCompletedHandler(
        [this]
        {
            uiView->
                setProgress(1.0);
        });
\end{verbatim}

Problems:

\texttt{uiView} may have been destroyed.

Callback thread may not be main.

\texttt{this} may have died.

UI state may race.

This is several bugs hidden in five lines.

\begin{center}\rule{0.5\linewidth}{0.5pt}\end{center}

\textbf{Safer architecture}

GPU callback produces plain data:

\begin{verbatim}
struct CompletedFrameInfo
{
    uint64_t Frame;
    double GpuMs;
};
\end{verbatim}

Then places it into thread-safe storage.

Main/UI side polls or receives main-thread dispatch.

Conceptually:

\begin{verbatim}
commandBuffer->addCompletedHandler(
    [weakState](...)
    {
        if (auto state =
                weakState.lock())
        {
            state->
                PushCompleted(...);
        }
    });
\end{verbatim}

Then UI:

\begin{verbatim}
DispatchQueue.main.async {
    consumeResult()
}
\end{verbatim}

Ownership and threading are explicit.

\begin{center}\rule{0.5\linewidth}{0.5pt}\end{center}

\textbf{Completion handlers and renderer destruction}

Maveb's renderer destructor drains its three frame-semaphore tokens.

Conceptually it waits until all submitted frames have completed before destruction continues.

That protects callbacks/resources referencing renderer state.

This is the Metal analogue of waiting/retiring in-flight frames before Vulkan teardown.

\begin{center}\rule{0.5\linewidth}{0.5pt}\end{center}

\textbf{Deferred cleanup in Maveb FrameContext}

Each frame context contains:

\begin{verbatim}
deferredCleanup_
\end{verbatim}

and:

\begin{verbatim}
deferUntilReusable(...)
\end{verbatim}

At:

\begin{verbatim}
beginFrame()
\end{verbatim}

it runs those cleanup functions.

Because frame slot is only reusable after semaphore completion discipline, this provides a simple frame-retirement mechanism.

Same idea as Vulkan deferred destruction.

Different synchronization primitive.

\begin{center}\rule{0.5\linewidth}{0.5pt}\end{center}

\textbf{This is the recurring cross-API pattern}

Vulkan:

\[
frameFence/timeline
\rightarrow
resource retirement.
\]

Metal:

\[
completion handler/frame semaphore/shared event
\rightarrow
resource retirement.
\]

The high-level engine concept is identical:

\[
\boxed{
\text{destroy only when last GPU use has completed}.
}
\]

This is precisely why the upcoming multi-backend abstraction must model \textbf{intent}, not API object names.

\begin{center}\rule{0.5\linewidth}{0.5pt}\end{center}

\textbf{Kairo Metal versus Maveb Metal}

The two backends represent two maturity points.

{\def\LTcaptype{none} % do not increment counter
\begin{longtable}[]{@{}
  >{\raggedright\arraybackslash}p{(\linewidth - 4\tabcolsep) * \real{0.3333}}
  >{\raggedright\arraybackslash}p{(\linewidth - 4\tabcolsep) * \real{0.3333}}
  >{\raggedright\arraybackslash}p{(\linewidth - 4\tabcolsep) * \real{0.3333}}@{}}
\toprule\noalign{}
\begin{minipage}[b]{\linewidth}\raggedright
Concern
\end{minipage} & \begin{minipage}[b]{\linewidth}\raggedright
Kairo Metal
\end{minipage} & \begin{minipage}[b]{\linewidth}\raggedright
Maveb Metal
\end{minipage} \\
\midrule\noalign{}
\endhead
\bottomrule\noalign{}
\endlastfoot
Device & default Metal device & supplied/retained device \\
Queue & one command queue & one labelled primary queue \\
Shader library & runtime source compilation & offline \texttt{.metallib} \\
Frame command buffer & one/frame & one complex command buffer/frame \\
Frames in flight & effectively simple/last-command model & 3 \\
Frame limiting & waits where needed & semaphore count 3 \\
Uploads & direct/shared buffers & per-frame 8 MiB shared arena \\
Readback & direct/shared + waits for simple operations & 256 KiB frame readback + async capture \\
Render targets & shared in simple backend & private \\
Transient heap & none in simple backend & 64 MiB private/frame \\
Hazard mode & default & tracked transient heap \\
Presentation & direct \texttt{CAMetalLayer} & \texttt{MTKView} drawable \\
Capture & waits on \texttt{LastCommand} & asynchronous blit + completion callback \\
GPU timing & not central & \texttt{GPUStartTime/GPUEndTime} \\
Labels & limited & extensive \\
UI bridge & renderer-neutral runtime & Objective-C++ + SwiftUI editor bridge \\
\end{longtable}
}

Neither codebase is ``wrong.''

Kairo prioritizes:

\[
\boxed{
\text{small comprehensible backend}
}
\]

while Maveb prioritizes:

\[
\boxed{
\text{overlapped multi-pass production behavior}.
}
\]

\begin{center}\rule{0.5\linewidth}{0.5pt}\end{center}

\textbf{A more production-oriented Kairo Metal evolution}

The most useful progression would be:

\[
\text{single LastCommand wait}
\]

to:

\[
\text{FrameContext[2 or 3]}.
\]

Each frame context gets:

\begin{verbatim}
struct MetalFrameContext
{
    id<MTLBuffer> Upload;
    size_t UploadCursor;

    std::vector<
        std::function<void()>>
        DeferredCleanup;
};
\end{verbatim}

At frame begin, wait only if rotating into a still-busy slot.

Then encode/commit.

Completion handler releases slot.

Exactly the model Maveb proves works.

\begin{center}\rule{0.5\linewidth}{0.5pt}\end{center}

\textbf{Kairo static meshes should eventually become private}

Current Metal mesh buffers are created with:

\[
StorageModeShared.
\]

For static geometry:

\[
CPU writes once,
GPU reads repeatedly.
\]

Better architecture:

\begin{verbatim}
shared upload buffer
    ↓
blit
    ↓
private vertex/index buffer
\end{verbatim}

Whether this materially improves Apple-silicon performance should be benchmarked, but it is the more general GPU-local architecture and mirrors Maveb's static upload path.

\begin{center}\rule{0.5\linewidth}{0.5pt}\end{center}

\textbf{Maveb's upload path demonstrates this}

Its static mesh upload creates destination and staging buffers.

Then:

\begin{verbatim}
blit->
    copyFromBuffer(
        staging,
        destination);
\end{verbatim}

Completion handler keeps staging object alive until copy finishes.

This is a subtle but important lifetime technique.

The lambda captures \texttt{staging}.

Therefore staging cannot be destroyed before GPU finishes copy.

Very elegant.

\begin{center}\rule{0.5\linewidth}{0.5pt}\end{center}

\textbf{Command-buffer retention of resources}

Metal generally retains resources referenced by encoded commands sufficiently for execution semantics, but your wrapper/application ownership still matters around buffers not directly retained in the same way, indirect pointers, callbacks and CPU-side state.

Do not treat command-buffer retention as an excuse to forget logical resource lifetime.

The engine still needs a coherent ownership model.

\begin{center}\rule{0.5\linewidth}{0.5pt}\end{center}

\textbf{Apple GPU thinking versus desktop immediate thinking}

The biggest practical mental shift is this:

Desktop-style intuition often optimizes around:

\[
\boxed{
\text{VRAM bandwidth + large caches + immediate framebuffer writes}.
}
\]

Apple GPU intuition should additionally prioritize:

\[
\boxed{
\text{tile locality}
}
\]

\[
\boxed{
\text{attachment lifetime}
}
\]

\[
\boxed{
\text{load/store elimination}
}
\]

\[
\boxed{
\text{unified-memory traffic}
}
\]

\[
\boxed{
\text{render-pass grouping}.
}
\]

This does not mean every Apple GPU is identical.

It means the tile-based architecture strongly affects good rendering structure.

\begin{center}\rule{0.5\linewidth}{0.5pt}\end{center}

\textbf{Why unnecessary pass splitting can hurt}

Imagine:

\begin{verbatim}
Pass A
    render color
    STORE to memory

Pass B
    LOAD color
    modify
    STORE

Pass C
    LOAD
    modify
    STORE
\end{verbatim}

Each pass boundary can force memory traffic.

If operations could legally remain in one rendering scope or use tile-local mechanisms, you may save enormous bandwidth.

This is one motivation behind render-pass/subpass/tile-aware graph optimization.

\begin{center}\rule{0.5\linewidth}{0.5pt}\end{center}

\textbf{Why over-fusing can also hurt}

One giant pass can:

\begin{itemize}
\tightlist
\item
  require too many attachments,
\item
  increase shader complexity,
\item
  force resources to remain live,
\item
  prevent async compute,
\item
  reduce modularity.
\end{itemize}

Again, the objective is not:

\begin{quote}
fewer passes at any cost.
\end{quote}

It is:

\[
\boxed{
\text{minimize unnecessary external memory traffic while preserving useful execution structure}.
}
\]

\begin{center}\rule{0.5\linewidth}{0.5pt}\end{center}

\textbf{Attachment lifetime equation}

For attachment \(R\):

If previous content is irrelevant:

\[
LoadCost(R)=0
\]

using Clear/DontCare semantics.

If result never needed after pass:

\[
StoreCost(R)=0.
\]

Ideal transient attachment:

\[
\boxed{
Load=DontCare/Clear
,\quad
Store=DontCare
}
\]

and memoryless where supported.

That resource can conceptually live entirely inside tile execution.

\begin{center}\rule{0.5\linewidth}{0.5pt}\end{center}

\textbf{Why temporal buffers cannot be transient}

TAA history:

\[
H_t
\]

must be consumed by frame:

\[
t+1.
\]

Therefore its lifetime crosses frame boundary.

It must:

\[
Store
\]

and remain allocated.

This is the opposite of memoryless.

The renderer must distinguish:

\[
\boxed{
\text{transient frame-local resource}
}
\]

from:

\[
\boxed{
\text{history/persistent resource}.
}
\]

\subsubsection{First-use completion --- lifetime analysis for a frame-local attachment and a history image}\label{05-engine-i-ch024-first-use-completion-lifetime-analysis-for-a-frame-local-attachment-and-a-history-image}

Here is the actual render-graph reasoning now, while transient and persistent storage first affect the result. Consider four operations: \textbf{geometry} writes a screen-space normal attachment \(N_k\); \textbf{lighting} reads \(N_k\) and writes a linear HDR image \(C_k\); \textbf{temporal accumulation} reads \(C_k\) and the preceding history \(H_{k-1}\) and writes new history \(H_k\); \textbf{presentation} reads \(H_k\). These are logical resource \emph{versions}. The graph edges are \[
\text{geometry}\to\text{lighting}\to\text{temporal}\to\text{present}
\] with another edge \(H_{k-1}\to\text{temporal}\) arriving from the preceding frame. An edge describes both the required execution ordering and the visibility of the producer's writes to the consumer, using the exact API's synchronization mechanisms. A memory allocation remaining alive is not proof that this happens-before relationship exists.

The graph can find a resource's logical \textbf{live interval} by recording the first writing pass and the last pass that must read it. The normal attachment \(N_k\) is needed from geometry through lighting but does not have to survive the frame when no later pass samples it. In a compatible tile-local rendering scope, and only where the platform/API supports the requested attachment usage, that can permit clear/dont-care load/store decisions or memoryless storage. The HDR image \(C_k\) must survive until the temporal pass reads it, so the compiler cannot discard it immediately after lighting. The history \(H_{k-1}\) must survive from its previous producer through \emph{this} temporal read, and \(H_k\) must survive until the following frame's temporal pass. Memoryless storage cannot fulfill that cross-frame read contract.

For an ordinary RGBA16F history, each pixel stores four 16-bit channels: \(4\times2=8\) bytes. At \(1920\times1080\) pixels one history uses \(1920(1080)(8)=16{,}588{,}800\) bytes, about \(15.82\) MiB, before alignment, tiling or additional metadata. Two such images alone account for about \(31.64\) MiB. The word \textbf{persistent} does not say they must always live in dedicated physical allocations, but it says that their contents cannot be thrown away while future work still needs them.

Two history versions are also not automatically sufficient for safe frames-in-flight. Frame \(k\) can read image A and write B. Frame \(k+1\) may read B and intend to write A, but A can be recycled only after the last asynchronous read that uses A has completed. A ring of physical images plus completion-aware retirement may be necessary. The distinction is between a \textbf{logical version} (the meaning of a frame's history) and a \textbf{physical allocation} (memory that can carry one version when no overlapping use forbids it). Conversely, allocating a separate texture for every named graph version may waste large amounts of memory if two transient versions have provably disjoint lifetimes and compatible usage.

\begin{verbatim}
geometry:
    WRITE N[k]                   # valid normal/color encoding, declared format
lighting:
    READ  N[k]                   # require producer-write visibility
    WRITE C[k]
temporal:
    READ  C[k]
    READ  H[k-1]                 # prior-frame content, not memoryless
    WRITE H[k]                   # new version, not an alias of a live input
present:
    READ H[k]

# The compiler can alias physical storage only if lifetime/schedule
# non-overlap, format/usage restrictions and synchronization all hold.
\end{verbatim}

One important qualification is that a linear list of passes does \textbf{not} automatically establish safe aliasing. With asynchronous compute or multiple command queues, two apparently non-overlapping entries in a textual schedule may actually execute concurrently. Aliasing is safe only after the compiler proves no overlapping physical access in the \emph{scheduled} dependency graph and establishes the required reuse barrier or completion. If a pass writes a resource and the next reads it, aliasing does not fix a missing visibility edge; if a pass reads a resource after its allocation has been reused, a correct barrier elsewhere does not restore the lost contents.

Temporal \textbf{validity} adds a second, semantic condition. Keeping \(H_{k-1}\) allocated and synchronized does not make its pixels meaningful after a camera teleport, resolution change, projection-convention change, material edit, estimator change or discontinuity in object identity. A history epoch should identify which camera/scene state produced the previous estimate; reset or deliberately reproject according to documented validity rules. A renderer that displays a stable but ghosted image can have perfectly correct GPU lifetime synchronization and still have incorrect history semantics.

At first use, therefore, a render-graph resource requires five answers: \textbf{who creates this version, who last consumes it, which operation makes writes visible, what completion permits storage reuse, and what semantic state makes its contents valid?} A test should intentionally delay one GPU submission and verify no buffer is recycled early; another should change the camera and verify history invalidation. Those are different tests. This is the complete reasoning needed here, not a placeholder for a later render-graph pass.

\begin{center}\rule{0.5\linewidth}{0.5pt}\end{center}

\textbf{Metal resource synchronization mental model}

You can now think in four levels.

Same encoder:

\[
\text{use barriers where shader memory dependency requires}.
\]

Different encoders, tracked resource:

\[
\text{Metal hazard tracking commonly orders required use}.
\]

Different encoders, untracked resource:

\[
\text{explicit fence/barrier semantics may be necessary}.
\]

Different queues/command streams:

\[
\text{events/shared events encode inter-stream dependency}.
\]

This hierarchy prevents treating all Metal synchronization as one thing.

\begin{center}\rule{0.5\linewidth}{0.5pt}\end{center}

\textbf{Metal/Vulkan naming trap}

Never build a cross-platform abstraction by matching names:

\begin{verbatim}
VkFence ↔ MTLFence
\end{verbatim}

because semantics are different.

Instead map concepts:

{\def\LTcaptype{none} % do not increment counter
\begin{longtable}[]{@{}
  >{\raggedright\arraybackslash}p{(\linewidth - 4\tabcolsep) * \real{0.3333}}
  >{\raggedright\arraybackslash}p{(\linewidth - 4\tabcolsep) * \real{0.3333}}
  >{\raggedright\arraybackslash}p{(\linewidth - 4\tabcolsep) * \real{0.3333}}@{}}
\toprule\noalign{}
\begin{minipage}[b]{\linewidth}\raggedright
Engine concept
\end{minipage} & \begin{minipage}[b]{\linewidth}\raggedright
Vulkan
\end{minipage} & \begin{minipage}[b]{\linewidth}\raggedright
Metal
\end{minipage} \\
\midrule\noalign{}
\endhead
\bottomrule\noalign{}
\endlastfoot
CPU knows GPU submission completed & \texttt{VkFence} / timeline value & completion handler / shared event value \\
GPU queue-to-queue timeline & timeline semaphore & shared event \\
encoder-local GPU dependency & barrier & memory barrier \\
resource dependency between phases & pipeline/image/buffer barrier & tracked hazard / fence / event depending scope \\
presentation image & swapchain image & drawable texture \\
\end{longtable}
}

This is the beginning of correct multi-backend thinking.

Pass 14 will formalize it.

\begin{center}\rule{0.5\linewidth}{0.5pt}\end{center}

\textbf{ARC versus GPU lifetime}

ARC answers:

\[
\boxed{
\text{Is the Objective-C object alive?}
}
\]

It does not answer:

\[
\boxed{
\text{Has GPU stopped using the object's underlying resource?}
}
\]

Those are different questions.

Likewise C++ \texttt{unique\_ptr} answers CPU ownership.

GPU completion remains another dimension.

Same lesson as Vulkan RAII.

\begin{center}\rule{0.5\linewidth}{0.5pt}\end{center}

\textbf{A robust bridge lifetime diagram}

\begin{verbatim}
SwiftUI AetherViewport
       ↓ owns
AetherViewportView
       ↓ owns
AetherViewportDelegate
       ↓ C++ unique_ptr owns
aether::metal::Renderer
       ↓ MetalPtr owns
MTL resources

GPU command buffer
       ↓ temporarily uses
MTL resources

completion
       ↓
frame slot becomes reusable
\end{verbatim}

Every arrow has a different ownership/synchronization meaning.

Understanding those arrows is much more valuable than memorizing ARC syntax.

\begin{center}\rule{0.5\linewidth}{0.5pt}\end{center}

\textbf{Swift closure lifetime}

A view property such as:

\begin{verbatim}
@property(copy)
void (^onEntityPicked)(...);
\end{verbatim}

owns a copied block.

Swift closure assigned to it can capture values/objects.

If it strongly captures the same object that strongly owns the view, you can form a retain cycle.

Therefore class-heavy bridges often use:

\begin{verbatim}
[weak self]
\end{verbatim}

when appropriate.

In Aether's current \texttt{NSViewRepresentable} closures, the bindings/value context differs from a typical view-controller retain cycle, but the general rule remains:

\[
\boxed{
\text{inspect closure capture graph explicitly}.
}
\]

\begin{center}\rule{0.5\linewidth}{0.5pt}\end{center}

\textbf{Thread ownership is part of API design}

Instead of documentation:

\begin{quote}
callback occurs later.
\end{quote}

Write:

\begin{quote}
callback may execute on a Metal completion thread; callback must be nonblocking and may not access AppKit. Results are marshalled to main thread.
\end{quote}

That is a real API contract.

Similarly:

\begin{quote}
\texttt{draw()} is invoked by MTKView frame lifecycle.
\end{quote}

\begin{quote}
\texttt{setCameraState()} must be called from UI/main integration layer.
\end{quote}

Threading assumptions should be first-class.

\begin{center}\rule{0.5\linewidth}{0.5pt}\end{center}

\textbf{Avoid holding locks in GPU completion handlers}

Completion handlers should be quick.

A callback that blocks on:

\[
heavy\ disk\ IO
\]

or a contended mutex can delay callback processing.

Good handler:

\begin{verbatim}
copy minimal metadata
move finished payload
atomic counter
signal semaphore
enqueue lightweight task
\end{verbatim}

Then larger work occurs elsewhere.

Maveb's handler mostly follows this pattern.

\begin{center}\rule{0.5\linewidth}{0.5pt}\end{center}

\textbf{Synchronous picking caveat in Maveb}

Some current picking paths create a command buffer, blit one pixel and call:

\begin{verbatim}
waitUntilCompleted();
\end{verbatim}

This is simple and correct.

But it stalls caller until GPU finishes.

For occasional editor clicks that may be acceptable.

For continuous hover picking at 60--120 Hz, convert to asynchronous frame readback just like Kairo's Vulkan pick path.

This is exactly the kind of design decision that depends on workload frequency.

\begin{center}\rule{0.5\linewidth}{0.5pt}\end{center}

\textbf{Kairo Metal's \texttt{WaitForLastCommand()}}

Kairo's destruction/pick/capture path waits for:

\[
LastCommand.
\]

Again:

\[
\boxed{
\text{simple and correct}
}
\]

but it serializes.

Its next maturity step is not ``remove the wait.''

It is:

\[
\boxed{
\text{replace global wait with explicit frame/resource retirement}.
}
\]

Never remove synchronization without adding an ownership model.

\begin{center}\rule{0.5\linewidth}{0.5pt}\end{center}

\textbf{The three Metal memory questions}

For every resource ask:

\[
\boxed{
\text{Does CPU need direct access?}
}
\]

If yes, Shared/Managed may be appropriate.

Then:

\[
\boxed{
\text{Does data need survive the render pass/frame?}
}
\]

If no, memoryless/transient strategy may apply.

Then:

\[
\boxed{
\text{Who synchronizes hazards?}
}
\]

Tracked mode or explicit synchronization.

These three questions often determine the correct Metal resource architecture.

\begin{center}\rule{0.5\linewidth}{0.5pt}\end{center}

\textbf{The three encoder questions}

Before adding an encoder ask:

\[
\boxed{
\text{Is this rasterization, compute, or transfer?}
}
\]

Then:

\[
\boxed{
\text{What resources does it read/write?}
}
\]

Then:

\[
\boxed{
\text{Can its work remain on-chip or does it force resource store/load?}
}
\]

This prevents accidental pass fragmentation.

\begin{center}\rule{0.5\linewidth}{0.5pt}\end{center}

\textbf{Practical Maveb optimization candidates revealed by M140}

Your current architecture already does many strong things:

\[
Private\ render\ targets,
\]

\[
3\ frames\ in\ flight,
\]

\[
Shared\ per-frame\ upload,
\]

\[
Shared\ readback,
\]

\[
Private\ transient\ heap,
\]

\[
Tracked\ hazards,
\]

\[
offline\ metallib,
\]

\[
labels,
\]

\[
async\ frame\ capture.
\]

The next questions are therefore not basic Metal correctness.

They are more advanced:

Can some strictly transient attachments become memoryless where supported?

Can graph lifetime analysis allocate actual temporary resources from the per-frame heap and alias them?

Can direct bindings become argument-buffer/resource-table architecture where material scale demands it?

Can synchronous pick commands become asynchronous if interaction frequency rises?

Can detailed counter sampling show whether Gaussian work is compute-, bandwidth- or tile-pressure-limited?

Those are mature optimization questions.

\begin{center}\rule{0.5\linewidth}{0.5pt}\end{center}

\textbf{PASS 13 RECONSTRUCTION TABLE}

{\def\LTcaptype{none} % do not increment counter
\begin{longtable}[]{@{}
  >{\raggedright\arraybackslash}p{(\linewidth - 2\tabcolsep) * \real{0.5000}}
  >{\raggedright\arraybackslash}p{(\linewidth - 2\tabcolsep) * \real{0.5000}}@{}}
\toprule\noalign{}
\begin{minipage}[b]{\linewidth}\raggedright
Topic
\end{minipage} & \begin{minipage}[b]{\linewidth}\raggedright
Reconstruction
\end{minipage} \\
\midrule\noalign{}
\endhead
\bottomrule\noalign{}
\endlastfoot
\texttt{MTLDevice} & GPU/capability/resource factory \\
\texttt{MTLCommandQueue} & ordered command-buffer submission stream \\
\texttt{MTLCommandBuffer} & one asynchronously executed GPU work package \\
Render encoder & rasterization commands \\
Compute encoder & general thread-grid kernels \\
Blit encoder & copies/mip generation/transfer work \\
\texttt{MTLBuffer} & linear byte resource \\
\texttt{MTLTexture} & structured image resource \\
\texttt{MTLSamplerState} & filtering/addressing/comparison \\
Render pipeline & compiled vertex+fragment and attachment-format/blend contract \\
Compute pipeline & compiled \texttt{kernel} state + execution properties \\
Metal library & container of compiled/runtime MSL functions \\
MSL vertex & \texttt{vertex} function producing \texttt{{[}{[}position{]}{]}} \\
MSL fragment & \texttt{fragment} function consuming interpolants/resources \\
MSL compute & \texttt{kernel} over thread grid \\
Shared & CPU+GPU accessible \\
Private & GPU-only optimized resource \\
Managed & CPU/GPU copies synchronized on applicable macOS configurations \\
Memoryless & transient tile-only attachment where supported \\
TBDR & bin geometry, rasterize per tile in on-chip memory \\
Load action & how attachment begins \\
Store action & whether tile result survives pass \\
Drawable & presentable texture instance \\
\texttt{CAMetalLayer} & drawable provider / Core Animation presentation layer \\
\texttt{presentDrawable} & schedule presentation through command buffer \\
Completion handler & asynchronous post-GPU callback \\
Frame semaphore & cap frames in flight \\
Upload arena & one shared per-frame buffer + aligned slices \\
Argument buffer & table of indirectly bound resources \\
Heap & pooled resource memory \\
Hazard tracking & automatic dependency tracking \\
Untracked resource & application owns synchronization \\
Memory barrier & shader/resource ordering within relevant execution scope \\
\texttt{MTLFence} & GPU encoder dependency primitive \\
Shared event & timeline-style inter-command/queue/CPU synchronization \\
GPU counters & hardware/timestamp measurements \\
Xcode capture & GPU workload inspection \\
Labels & semantic names for GPU tools \\
ARC & Objective-C reference-count automation \\
\texttt{.mm} & Objective-C++ translation unit \\
C++ bridge ownership & RAII and ARC coexist \\
Error bridge & C++ Result → NSError/BOOL/nil/etc. \\
SwiftUI role & declarative application/editor state \\
Renderer role & GPU resources/execution \\
Main-thread rule & UI work stays on UI thread \\
Renderer callback rule & thread and ownership contract must be explicit \\
\end{longtable}
}

\begin{center}\rule{0.5\linewidth}{0.5pt}\end{center}

\textbf{Blank-page mastery gate}

Reconstruct a Metal frame from nothing.

Start with:

\[
MTLDevice
\]

then create:

\[
MTLCommandQueue.
\]

At frame \(N\), acquire one frame context:

\[
F_{N\bmod3}.
\]

Explain why the frame semaphore must guarantee that \(F\)'s previous GPU use has completed before resetting:

\[
uploadCursor=0.
\]

Then create:

\[
MTLCommandBuffer.
\]

Encode a shadow render encoder.

End it.

Encode a compute encoder.

End it.

Encode the scene render encoder.

End it.

Encode TAA/bloom.

Acquire drawable.

Encode final presentation pass.

Schedule:

\[
presentDrawable.
\]

Add completion handler.

Commit.

Then explain why the CPU can immediately begin preparing a future frame without waiting for current GPU completion.

From there, derive why this resource:

\begin{verbatim}
CPU updates every frame
small
GPU reads
\end{verbatim}

should usually be:

\[
Shared.
\]

Why this:

\begin{verbatim}
large HDR render target
CPU never directly reads it
\end{verbatim}

should usually be:

\[
Private.
\]

Why this:

\begin{verbatim}
depth used only during one tile-local pass
never sampled later
\end{verbatim}

can potentially be:

\[
Memoryless.
\]

Then explain exactly why:

\[
LoadActionClear
\]

can save external memory bandwidth relative to \texttt{Load}, and why:

\[
StoreActionDontCare
\]

can eliminate writeback.

If you cannot connect load/store actions to Apple's tile memory, you do not yet understand Metal performance.

Then compare:

\[
MTLFence
\]

against:

\[
VkFence.
\]

Explain why mapping them by name would be wrong.

Then compare:

\[
MTLSharedEvent(value=N)
\]

with:

\[
VkTimelineSemaphore(value=N).
\]

Then reconstruct Maveb's \texttt{FrameContext}:

\[
8\text{ MiB Shared Upload}
+
256\text{ KiB Shared Readback}
+
64\text{ MiB Private Tracked Heap}.
\]

Explain why:

\[
offset=
(value+a-1)\&
\sim(a-1)
\]

requires \(a\) to be a power of two.

Then explain why:

\begin{verbatim}
frame.beginFrame();
\end{verbatim}

is safe only after that frame slot has retired.

Then reconstruct the ownership boundary:

\[
SwiftUI
\rightarrow
AetherViewport
\rightarrow
AetherViewportView
\rightarrow
AetherViewportDelegate
\rightarrow
std::unique_ptr<Renderer>
\rightarrow
MetalPtr<MTLResource>.
\]

Explain where ARC is responsible.

Explain where C++ RAII is responsible.

Explain why neither one by itself proves GPU completion.

Then explain why this bridge cast:

\begin{verbatim}
(__bridge void*)device
\end{verbatim}

does not transfer ownership, and why the persistent C++ renderer must retain the device if it stores it.

Then explain why a Metal completion handler must not blindly mutate SwiftUI state.

Correct chain is:

\[
GPU\ completion
\rightarrow
thread\text{-}safe\ result
\rightarrow
DispatchQueue.main
\rightarrow
UI.
\]

Finally compare your two implementations.

Kairo Metal currently chooses:

\[
\boxed{
\text{shared resources + simple last-command waiting}
}
\]

for clarity.

Maveb chooses:

\[
\boxed{
\text{private GPU resources + triple buffering + per-frame arenas + async completion}
}
\]

for throughput.

The difference is not API knowledge.

It is \textbf{lifetime and scheduling architecture}.

Once you can see that, Metal stops looking like a collection of Objective-C calls and becomes a coherent execution machine:

\[
\boxed{
\text{resources}
+
\text{encoder phases}
+
\text{tile locality}
+
\text{frame ownership}
+
\text{async lifetime}
}.
\]

This is where Phase 24 becomes genuinely engine-architectural: \textbf{backend-neutral handles and resource descriptions without lowest-common-denominator design; feature/capability contracts and backend escape hatches; then the complete render-graph system---resources, passes, reads/writes, dependency inference, cycle detection, topological scheduling, lifetime intervals, transient allocation and aliasing, barrier/state inference, cross-queue dependencies, pass culling, imported/history/swapchain resources and graph debugging---followed by Metal↔Vulkan conversion boundaries, portable mesh/material/scene extraction, private native handles, shader/reflection/layout strategy, coordinate/depth convention adaptation, backend visual agreement, and designing a multi-backend renderer that shares engine intent without pretending Vulkan and Metal are identical.}

Graph connections: \hyperref[appendix-f662b5b1773b]{GPU workgroups and synchronization} and \href{../04-Light-Transport-and-Image-Synthesis/main.pdf\#nameddest=appendix-f2ee4d086415}{Prefix scan and stream compaction (Light Transport and Image Synthesis, Reference notes: Prefix scan and stream compaction)}. The shared spatial/order/lifetime visual is \href{../07-Vision-XR-and-Playable-Worlds/main.pdf\#nameddest=04-visual-computing-ii-ch037-figure-er1---footprint-compositing-order-and-storage-lifetime}{Figure ER1 (Vision, XR, and Playable Worlds, Vision, XR, and Playable Worlds, Chapter 6)}.

\subsection{Distinguish host allocation callbacks from device-memory allocation}\label{05-engine-i-ch024-distinguish-host-allocation-callbacks-from-device-memory-allocation}

There are two allocator conversations inside a Vulkan renderer. An application allocates and binds device memory for buffers and images. Separately, Vulkan implementation code may need host memory for its own internal bookkeeping. \texttt{VkAllocationCallbacks} lets the application supply functions for that second purpose. A host allocation callback is not a replacement for \texttt{vkAllocateMemory}, and a counter attached to that callback is not a measurement of all resident graphics memory.

The allocation function receives a requested byte count and an address alignment. \textbf{Size} answers how many usable bytes must exist. \textbf{Alignment} answers which addresses are permitted for the returned pointer. A request for 37 bytes at 16-byte alignment does not require a 64-byte logical size: the returned address must be divisible by 16 and the block must contain at least 37 usable bytes. An allocator may round its internal consumption upward, but that is its storage policy. In particular, the power-of-two requirement belongs to alignment. An extracted caption claim that assigns it to the byte count must be corrected; the original subtitle text explicitly refers to alignment.

Reallocation asks to change a block while retaining the shared prefix of its contents. If an old 37-byte block becomes a 70-byte block, its first 37 bytes must remain. If it shrinks to 20 bytes, its first 20 remain. A replacement strategy can allocate the new block, copy the smaller byte count, then free the old block. It must not free the old block before it knows replacement succeeded. Otherwise failure destroys the caller's only valid copy.

Three boundary cases deserve explicit branches. A null original pointer requests allocation behavior. A zero requested size requests free behavior. When a non-null original block cannot be reallocated, the application must leave the old allocation intact. The requested alignment for reallocation of an existing block must agree with the alignment originally used. Free must also accept a null pointer safely. These are callback semantics, rather than a reason to guess at an apparent documentation typo. They are defined in the \href{https://docs.vulkan.org/spec/latest/chapters/memory.html}{Khronos allocation-callback specification}.

A useful trace begins with a 37-byte block containing an identifier and payload. A request for 70 bytes fails: the returned replacement is null, the original pointer still owns the same 37 bytes, and the caller can retry or report failure. A later successful request returns a valid aligned block with the prefix intact; only then does ownership move and the original block cease to be valid. Finally a zero-size request releases that allocation through the specified free behavior. An adapter must preserve these cases even if its underlying allocator uses a different convention.

The source episode also distinguishes internal-allocation notifications from application-supplied allocation/free callbacks. A notification reports an event for accounting; it does not hand the engine a new memory block to destroy. Mixing those paths can double-free memory or make diagnostics look more complete than they are. Keep each counter's measured scope beside its display: hooked host requests, application device allocations, driver estimates, and resident hardware usage answer different questions.

This clarification resolves the uncertainty noted in \href{../05-Spatial-Worlds-Engine-State-and-Simulation/main.pdf\#nameddest=05-engine-i-ch014-bridge-vulkan-allocation-callbacks-to-an-engine-allocator}{the engine allocator walkthrough (Spatial Worlds, Engine State, and Simulation, Spatial Worlds, Engine State, and Simulation, Chapter 14)}. It checks the interface contract against primary documentation; it does not certify the episode's implementation or every allocation path in Kohi.

\section{A generated binding is still an API boundary}\label{05-engine-i-ch024-a-generated-binding-is-still-an-api-boundary}

A code generator can remove repetitive wrapper work, but it does not erase the difference between languages. The Sparky branch in this devlog experiments with exposing selected C++ engine operations through C++/CLI. C++/CLI is a C++ extension for working with the .NET managed runtime while interoperating with native C++ code. The language extension is not a language-neutral ABI. Its wrapper calls the native API and converts values at the boundary---for example, a managed string must be converted into the representation expected by the native C++ method. The episode also describes limiting exposure to public or protected operations and says that the branch is a learning project with substantial work remaining. Treat this as a prototype snapshot rather than a production generator or complete C\# integration (Sparky Engine Dev Log \#12, 00:29:20--00:30:03; source \texttt{src-a539341c398d4c55cdb1}, unit \texttt{u-486c51f3035705147bf6ac71}, with the branch status in \texttt{u-461e47c9ea075d3a77cdba67}).

This small example exposes the actual generator contract. The native API has types, ownership rules, error behavior, and lifetimes. The target language has its own strings, arrays, exceptions, and object model. A wrapper must decide which native operations are visible, how each type is represented, who owns the converted value, how errors cross the boundary, and whether callbacks can safely outlive the wrapper object. Generating a method signature is easy compared with preserving those semantics.

A reliable binding workflow starts from a narrow API surface and an explicit mapping table. For each supported native type, define its target-language form and conversion, including ownership and nullability. For each function, identify which side allocates and releases memory, how errors are reported, and whether references are borrowed or retained. Generate a small wrapper, compile it, and exercise one call in both directions where callbacks are involved. Keep generated output reproducible from the declared source and tool version; inspect representative generated code during development so an accidental mapping does not become hundreds of opaque failures.

The devlog's filesystem section offers a parallel lesson: a virtual path such as \texttt{/models/hero} decouples an asset reference from a workstation's physical directory, while ordered mounts determine where a request resolves. Both mechanisms insert a stable vocabulary between a client and an implementation. The vocabulary is useful only when its resolution and ownership rules are explicit. The source says packed-asset support is future work, so no archive format or package performance is implied here (00:04:01--00:05:40, unit \texttt{u-700ebba84bd2d2a5f5cabfeb}, and future-plan point \texttt{p-000011}).

A generator does not bypass C++ access control: protected native operations still need a valid class/access context and a deliberate exposure policy. A managed string conversion must specify encoding and temporary-buffer lifetime, not merely make the types compile. A generated signature is evidence of a proposed interface; a compiled wrapper, a successful call and a correct ownership/error path are separate verification steps.

\subsection{Choose the language boundary before choosing a generator}\label{05-engine-i-ch024-choose-the-language-boundary-before-choosing-a-generator}

\textbf{Managed code} runs with a runtime's type and execution model; \textbf{native code} follows its compiled platform's binary interface and its own ownership contracts. C++/CLI can combine managed and native operations in a .NET bridge. That does not make a native pointer a garbage-collected owner or define how long a borrowed buffer remains valid. Microsoft documents mixed managed/native interoperability and describes C++/CLI as a Windows bridge retained primarily for existing systems. The old Sparky experiment is useful for understanding the boundary, rather than as a default architecture for every new engine. \href{https://learn.microsoft.com/en-us/cpp/dotnet/native-and-dotnet-interoperability}{Microsoft native/.NET interoperability}, \href{https://learn.microsoft.com/en-us/cpp/dotnet/}{C++/CLI scope}.

For a cross-language engine interface, another design exposes a narrow C-compatible boundary: create an opaque engine handle, perform operations through functions, return explicit status/results, and destroy the handle through its owning module. Keep C++ object layouts, exceptions and allocator assumptions behind that boundary. This does not make one binary portable across all operating systems; each target still has a platform ABI. It does make the supported calling surface easier to specify and bind. Microsoft's .NET guidance recommends a C-style boundary for ordinary C++ interoperability because C++ has no single consistent ABI across its supported compilers and platforms. Choose the actual runtime, platforms, callback needs and ownership model first; then select tooling that can preserve that contract. \href{https://learn.microsoft.com/dotnet/standard/native-interop/abi-support}{Microsoft native interoperability ABI guidance}.

\subsection{A visual description does not identify a shader operator}\label{05-engine-i-ch024-a-visual-description-does-not-identify-a-shader-operator}

The 2D-API narration describes leaving a red control at one and lowering green and blue controls to 0.9. It calls the effect ``subtracting'' those channels ``in a way.'' That language describes the visible suppression of green and blue; it does not identify the arithmetic operator in an uninspected shader. A standard multiplicative tint is fully compatible with the explanation:

\[
C_{out}=C_{texture}\odot(1,0.9,0.9),
\]

where \(\odot\) means component-wise multiplication in the declared shader working space. The red component is unchanged and the others are reduced. White factors \((1,1,1)\) are neutral for multiplication. A literal subtraction of those control values would have a very different neutral point, while subtraction of a separately computed offset could again look similar for a particular image.

To distinguish models, use known input texels and inspect the shader or measured output. For a channel value of \(0.5\), multiplying by \(0.9\) yields \(0.45\), whereas subtracting an offset of \(0.1\) yields \(0.4\). A saturated white input may hide that difference. This small witness explains why visual similarity and casual words are weak evidence about implementation. It does not prove what the historical episode's shader executed. Preserve the mathematical teaching model and label the exact episode operation unverified until its code or a controlled test establishes it.

\section{A platform escape hatch must also restore process-wide state}\label{05-engine-i-ch024-a-platform-escape-hatch-must-also-restore-process-wide-state}

The Linux-windowing episode gives a concrete example of why a backend boundary needs to include side effects, not only function signatures. The Kohi \#005 Linux platform episode uses Xlib and XCB together for an X11 window. It calls Xlib to disable key-repeat events because repeated key presses would complicate the engine's input handling. The presenter explicitly warns that the call is global to the X server/session rather than scoped to this window. If the process exits without turning repeat back on, the user may continue to experience changed keyboard behavior outside the game. The same episode therefore adds an explicit shutdown action to re-enable repeat before destroying its window. (Kohi \#005, source \texttt{src-994a368b666ebd1bf250}, units \texttt{u-81fe6c39aad8fdcf8055e768}, \texttt{u-f4883a4974cbd13730aea658}.)

The transferable engineering lesson is to model a platform call as a resource or state transition with a matching cleanup obligation. Record: what state is changed, who else can observe it, how the change is scoped, and what cleanup runs on normal close and failure. A platform adapter that exposes only \texttt{disable\_key\_repeat()} hides the important contract. A safer boundary either avoids changing global behavior, uses a window-scoped mechanism where available, or makes acquisition and restoration explicit and testable. If initialization fails after changing the setting, rollback still needs to occur. This is proposed API design based on the tutorial's warning, not a claim about all X11 input systems.

The episode's choice to use both Xlib and XCB is also a reminder to distinguish protocol concepts from bindings and libraries. It demonstrates one historically specific implementation path; its remarks about documentation and deprecation are the speaker's assessment at recording time, not current Linux platform guidance. A modern engine should first identify the presentation and input contracts it needs, then select and verify supported platform APIs for its target environments. The supplied source does not establish present-day X11 recommendations or behavior under Wayland.

\subsection{Restore the state that existed before the adapter changed it}\label{05-engine-i-ch024-restore-the-state-that-existed-before-the-adapter-changed-it}

Unconditionally enabling repeat on close is not the same as restoring the user's previous setting. Xlib exposes the current keyboard controls, including global auto-repeat mode, through \texttt{XGetKeyboardControl}; its global repeat calls affect the keyboard for the selected display. Record the initial state before a deliberate global change and restore that state on the supported cleanup paths. \href{https://www.x.org/releases/X11R7.6/doc/libX11/specs/libX11/libX11.html}{X.Org Xlib keyboard controls}.

A destructor or rollback handler can organize normal cleanup, but it cannot guarantee restoration after every abrupt process termination. Global state can also be changed by another client while the application runs. These limitations strengthen the case for an input design that avoids changing shared session preferences. If a platform workaround must change them, its scope, rollback and competing-user assumptions belong in the adapter contract alongside its function signatures. This is a design qualification of the historical tutorial, not a native X11 test on this machine.

\subsection{Range count, byte capacity and API guarantees are different quantities}\label{05-engine-i-ch024-range-count-byte-capacity-and-api-guarantees-are-different-quantities}

An array containing 32 \texttt{VkPushConstantRange} records stores 32 descriptions in host memory. It does not create 128 bytes of push-constant capacity on a device. Each record describes a byte offset, a byte size and the shader stages that use that range. Counting records and counting bytes answer different questions.

For a declared range, its end offset must fit the queried \texttt{maxPushConstantsSize}. For example, offset 64 and size 96 end at byte 160: this fails a 128-byte device limit and passes that particular extent check against 256 bytes. Alignment, stage flags and the pipeline's interface still have their own requirements. Passing one byte-count check does not validate an entire pipeline layout. \href{https://docs.vulkan.org/spec/latest/chapters/limits.html}{Khronos physical-device limits}.

The name \texttt{maxPushConstantsSize} reports this device's supported maximum; the specification separately defines a minimum capability that an implementation must provide. Vulkan 1.0 through 1.3 guarantee at least 128 bytes, while Vulkan 1.4 raises that required minimum to 256 bytes. A device can report more than its baseline minimum. Query the actual property and state the application's supported baseline before designing a block layout. A tutorial's 128-byte example is a useful portable baseline for its version, not a universal maximum for all Vulkan devices. \href{https://docs.vulkan.org/guide/latest/versions.html}{Khronos versions and limit changes}.

This distinction applies far beyond push constants: a wrapper's array length, an allocation's host byte size, an API version's guaranteed floor and a device's advertised limit are four separate facts. Keep their units and owners beside their names so a local implementation constant cannot quietly become a claimed hardware rule.

\begin{center}\rule{0.5\linewidth}{0.5pt}\end{center}

\hyperref[volume-reading-map]{Reading map} · \hyperref[05-engine-i-ch023-execute-and-measure-a-metal-frame]{Previous chapter} · \hyperref[06-engine-ii-ch025-renderer-evolution-through-resource-ownership]{Next chapter}
